\CatchFileBetweenTags{\BoundaryStorageVal}{constants.tex}{BoundaryStorageVal}
\CatchFileBetweenTags{\GaugeLoadVal}{constants.tex}{GaugeLoadVal}
\CatchFileBetweenTags{\TotalBandwidthVal}{constants.tex}{TotalBandwidthVal}

\section{Synthesis: Global Solvency of the Nested Architecture}
\label{sec:Synthesis}

In \cref{sec:nested_persistence}, we established that any persistent system whose internal complexity exceeds a single control loop must partition its bandwidth according to three strict structural rules. Having derived the gauge, scalar, and gravitational sectors of the vacuum independently, we must now demonstrate that they form a mathematically closed, interdependent network. For the $E_8$-Persistence theory to be valid, it must be \textbf{globally solvent}: the sum of the partitioned subsystems must perfectly exhaust the substrate, the minimal structural remainders must be conserved, and all nested boundaries must seamlessly connect.

\subsection{Validation of Rule 1: Capacity Partitioning (The \texorpdfstring{$\nu=16$}{nu=16} Limit)}
Under Rule 1, partitioning is a strict zero-sum allocation of the hardware capacity. The total allocated bandwidth must exactly match the fundamental chiral limit ($\nu = 16$) without overflow or dead space. 

By summing the derived capacity loads for the Bulk (Gravity), the Boundary (Topology), and the Gauge (Interaction), we verify the global conservation of information:
\begin{equation}
\underbrace{B_{res}}_{\text{Gravity (Bulk)}} + \underbrace{\frac{\chi}{\sigma-\chi}}_{\text{Matter (Boundary)}} + \underbrace{\alpha}_{\text{Synchronization (Gauge)}} \equiv \nu
\end{equation}

Substituting the derived geometric values from the previous sections:
\begin{equation}
\ResidualCapVal + \BoundaryStorageVal + \GaugeLoadVal = \TotalBandwidthVal \quad (\text{Exact})
\end{equation}

\textit{Note on Circularity:} The sum above is algebraically exact by the definition of $B_{res}$ as the partition remainder; it cannot fail. The non-trivial claim is cross-sector: the same $B_{res} \approx 15.326$ independently governs the gravitational depth in Section~\ref{sec:geometric_depth}, where it couples with the Beer-Lambert attenuation ($\alpha^{\Delta/2}$) to reproduce both $\alpha_G$ and $M_P$ at the correct experimental values. A numerically different remainder --- arising from any other partition scheme --- would yield incorrect gravitational predictions. The solvency check is therefore a consistency test of the partition structure, not a verification of the sum.

\subsection{Validation of Rule 2: Reversible Encoding (The Landauer Floor)}
Under Rule 2, a persistent system cannot simply erase geometric mismatches or fractional remainders, as this would incur an infinite dissipative heat penalty (Landauer's Principle). It must strictly conserve this minimal geometric resolution as structural metadata. 

The vacuum successfully satisfies this by establishing a strict, non-zero resolution floor: the \textbf{Margin Impedance} ($Z_{MAR}$). As derived in Section \ref{subsec:margin_electron}, the geometric dilution of a single unit of information across the lattice's phase-space establishes the electron mass ($m_e$). The architecture does not permit a mass of zero for charged boundaries; the electron represents the exact Landauer Limit of the substrate, the minimum viable topological defect that can be reversibly encoded without dissolving into thermal noise. Rule 1 sets the capacity ceiling ($\nu=16$); Rule 2 successfully sets the structural floor.

\subsection{Validation of Rule 3: Nested Impedance Matching}
Under Rule 3, when a system partitions into nested layers, information crossing the boundaries between them must be translated with minimal dissipative loss. This requires the internal processing resistance (Impedance) of the regulatory subsystems to perfectly match the geometric aperture of the signals they govern.

As derived in their respective sections, the vacuum successfully satisfies this requirement at both extremes of the substrate's geometry:
\begin{itemize}
    \item \textbf{The Surface Match (System 5):} The local Higgs regulator structurally matches the aperture of the force it mediates ($A_{weak} \approx 5.965$). Its internal impedance ($Z_H \approx 5.964$) minimizes entropic drag at the $\chi=2$ topological boundary, perfectly matching the vacuum resonance to the weak channel capacity.
    \item \textbf{The Core Match (System 6):} The global Gravitational regulator saturates the lattice at the core scale ($r=\Delta/2$). The Planck Mass ($M_P$) is the exact, unique scale required to normalize the bulk geometric attenuation such that the core-to-surface impedance achieves exact unity ($Z_G \equiv 1$).
\end{itemize}

\subsection{Final Verdict: Global Solvency}
The Effective Field Architecture of the vacuum is globally solvent. Its $\nu = 16$ capacity budget is exhausted across the gauge, matter, and gravitational sectors without overflow or dead space. Its structural floor is set by the Landauer limit of the electron, and its nested boundaries are impedance-matched from the topological surface to the geometric core.

The fundamental constants of the Standard Model, the dimensionless couplings (Tier 1), the massive energy scales (Tier 2), and the time-asymmetric causal flow ($J_{CKM}$), are not independent, arbitrary parameters. They are the unique, strictly constrained solutions to a global impedance-matching problem. 

The physical universe persists because its substrate transfers information from the lattice core ($r = \Delta/2$) to the surface ($r = 0$) with minimal dissipative loss, reversibly encoding its boundary states (matter) while flawlessly exhausting its finite $\nu=16$ capacity budget.